\chapter{Blood Transfusion}

\section*{Overview}
Blood transfusion is the intravenous administration of whole blood or a blood component, such as packed red blood cells, platelets, plasma, or cryoprecipitate, to replace deficient blood elements. It is a common supportive therapy for anemia, hemorrhage, and coagulation disorders.

\section*{Alternative Names}
Red blood cell transfusion, packed red cell transfusion, platelet transfusion, fresh frozen plasma transfusion, cryoprecipitate transfusion

\section*{Principle}
Transfusion replaces circulating blood volume and deficient cellular or protein components, restoring oxygen-carrying capacity, hemostasis, and oncotic pressure. Compatibility is ensured by ABO and Rh typing with crossmatching, so that transfused cells and proteins are not destroyed by recipient antibodies. Hemovigilance systems monitor for transfusion reactions and infectious complications.

\section*{History}
The first human blood transfusion was recorded in the 17th century, and safe, compatible transfusion became practical after Karl Landsteiner described the ABO blood groups in 1900. Component therapy and modern blood banking evolved during the 20th century.

\section*{Why This Test or Procedure Is Ordered}
Ordered for symptomatic anemia, acute hemorrhage with hemodynamic instability, severe thrombocytopenia or platelet dysfunction with bleeding, coagulation factor deficiencies, and exchange transfusion in severe neonatal jaundice or sickle cell crisis. It is also used to support patients undergoing surgery, chemotherapy, or bone marrow transplantation.

\section*{Is This a Primary or Secondary Test or Procedure}
A primary therapeutic procedure when specific blood components are required to treat acute or life-threatening deficiency.

\section*{How This Test or Procedure Is Performed}
Blood is collected into bags containing anticoagulant-preservative solutions and processed into components. After type and crossmatch verification, the unit is administered through a dedicated intravenous line with a filter, usually over one to four hours. Vital signs are monitored before, during, and shortly after the transfusion, which is stopped immediately if a reaction is suspected.

\section*{What Preparation Is Needed}
Informed consent, ABO and Rh typing, antibody screening, and crossmatch of donor and recipient blood. Premedication with antipyretics is sometimes considered for patients with prior febrile reactions.

\section*{Contraindications and Precautions}
Transfusion is avoided when alternative treatments are effective, such as iron, folate, or vitamin B12 therapy for nutritional anemia, and in patients with antibody-mediated incompatibility. Precautions include strict patient identification to prevent mistransfusion, leukoreduction and irradiation for immunocompromised recipients, and caution in heart failure or volume overload states.

\section*{Risks}
Febrile nonhemolytic and allergic reactions are most common; more serious risks include acute and delayed hemolytic reactions, transfusion-related acute lung injury (TRALI), transfusion-associated circulatory overload (TACO), and transmission of infectious agents.

\section*{What Happens After the Test or Procedure}
The patient is observed for several hours for signs of a transfusion reaction, and follow-up complete blood counts are checked to document response. Recurring needs may prompt a search for the cause of ongoing blood loss or marrow failure.

\section*{Understanding Your Results}
The goal is clinical improvement, reflected by a rising hemoglobin, platelet count, or correction of coagulopathy to a target range. Results that fail to improve suggest ongoing losses, dilution, or refractoriness to transfusion.

\section*{Normal Results}
Hemoglobin concentration, platelet count, and coagulation parameters improve toward age- and sex-appropriate reference ranges after transfusion.

\section*{Follow-up Tests and Procedures}
Repeat blood counts and coagulation studies guide further transfusions. Patients with repeated reactions or platelet refractoriness may require antibody testing and HLA-matched components.

\section*{References}
\begin{enumerate}
  \item MedlinePlus. Blood transfusion. U.S. National Library of Medicine; 2025.
  \item Current Medical Diagnosis and Treatment. McGraw-Hill; 2025.
\end{enumerate}

\clearpage
